\subsection{Cross-Dataset Variance Analysis}

To assess algorithm stability across different data characteristics, we computed the variance of grades received by each algorithm-feature pair across all 80 datasets. The variance was calculated on a 0--4 scale (F=0, D=1, C=2, B=3, A=4), where lower variance indicates consistent performance across diverse time series types. Table~\ref{tab:variance} presents the complete variance matrix, with darker shading highlighting algorithm-feature combinations that exhibit greater sensitivity to data characteristics.

Algorithm-feature pairs are distributed across three distinct variance categories: \textbf{consistent} (variance $<0.5$, indicating stable performance across datasets), \textbf{moderately data-dependent} (variance $0.5$--$1.5$, showing 1--2 letter grade fluctuations), and \textbf{highly data-dependent} (variance $>1.5$, exhibiting fluctuations exceeding 2 letter grades). Variance values range from 0.01 (Douglas-Peucker on roughness, slope $\ell_1$, and level $\ell_1$) to 3.39 (Min Filter on extrema $W_\infty$).

Several distinct patterns emerged when analyzing variance by algorithm family:

      c
\paragraph{Transformers} display the widest range of variance patterns, spanning all three categories depending on their algorithmic approach. IIR frequency filters (Butterworth, Chebyshev, Elliptical) demonstrate remarkable consistency with variance predominantly below 0.5 across most metrics, though showing moderate variance ($1.0$--$1.4$) for regime detection. Convolutional smoothers (Gaussian Filter, Mean Filter, Savitzky-Golay) exhibit moderate variance in the $0.5$--$1.5$ range for most metrics, with Gaussian Filter ranging from 1.12 to 1.46 and Mean Filter from 1.12 to 1.51. Savitzky-Golay shows similar moderate variance but reaches high variance ($>1.5$) for regimes (1.91) and spikes/dips $W_1$ (1.56), indicating greater sensitivity to abrupt changes. Extreme-value filters (Min Filter, Max Filter, Median Filter) exhibit the highest variance overall, with Min Filter reaching 3.39 for extrema $W_\infty$, 3.04 for extrema $W_1$, and 2.69 for spikes/dips $W_1$, while Median Filter shows 2.76 for spikes/dips $W_1$ and 2.46 for extrema $W_1$. FFT Cutoff Filter demonstrates mixed behavior with high variance for noise $\ell_\infty$ (1.92) but consistency elsewhere.

\paragraph{Aggregators} occupy the moderate variance category. ASAP shows variance predominantly in the $1.0$--$1.5$ range across most metrics, reaching 1.83 for periodicity amplitude detection. PAA demonstrates slightly lower variance ($0.5$--$1.0$) for most features, suggesting more stable behavior than ASAP's adaptive windowing approach.

These variance patterns reveal a critical distinction between \textbf{consistent excellence}, \textbf{robust general-purpose performance}, and \textbf{data-dependent specialization}. Low-variance algorithms with low grades (e.g., Douglas-Peucker) are consistently inadequate. High-variance algorithms (e.g., Min/Max/Median filters) show dramatic performance swings, excelling on specific datasets but failing on others. Most notably, the top-performing algorithms (Gaussian Filter with GPA 3.70, Mean Filter with GPA 3.675, Savitzky-Golay with GPA 3.59) exhibit moderate variance ($0.5$--$1.5$), indicating \textbf{robust general-purpose smoothers} that reliably achieve high grades (A or B) across most datasets while still showing 1--2 letter grade fluctuations based on data characteristics. This moderate variance does not indicate unreliability, but rather appropriate adaptability: these algorithms consistently perform well without being uniformly optimal, making them safe default choices when data characteristics are unknown.
